\chapter*{Conventions}
\addcontentsline{toc}{chapter}{Conventions}

\paragraph{Numbers come from a ledger.} Every number derived from data is a field of a claim in the ledger of
\cref{app:claims-ledger}, generated at build time from the artefact it was computed from. A count is printed exactly;
other values are rounded to a few leading digits, a precision the claim may set, and a share is printed
either as a fraction or, where the text says so, as a per cent. Constants of the analysis, such as thresholds, interval
levels and the dates that bound the dataset, come from the frozen protocol and are printed the same way. A number that
the build cannot resolve prints as a red marker and fails the test suite, so a printed copy never contains a number
that was typed by hand.

\paragraph{Intervals.} Unless a claim says otherwise, an interval is at level $\proto{intervals.level}$ and is written
$[\,\text{lower}, \text{upper}\,]$. Its method is recorded in the ledger and chosen as \cref{ch:evidence-model}
describes: over the independent unit when rows are nested in customers, agents, branches or campaigns, over blocks of
days when the series is serially dependent, and under independence of rows only when rows are independent. An interval
that the notebooks computed treating dependent rows as independent is marked \naivemark{}; it is likely too narrow and
is reported as found, not used for a decision.

\paragraph{Status.} Each claim has a status that limits its use.
\begin{description}[style=nextline]
  \item[Exploratory] Found by the exploratory notebooks or recomputed by this report; it describes the data and may
  guide further work.
  \item[Registered] Re-estimated under a specification frozen before it ran; only a registered claim decides a
  readiness gate.
  \item[Positive control] Measured on data with a planted answer; it shows what a method can detect and says nothing
  about the bank.
\end{description}

\paragraph{Tests and discoveries.} A sentence that calls a result statistically detectable cites its
false-discovery adjusted value $q$, computed over a family of tests that the protocol declares, and a result is a
discovery when $q \le \proto{gates.alpha}$. A p-value cited without $q$ belongs to a single declared test.

\paragraph{Tables and figures.} A table rendered from a stored result prints, in its caption, the path of that result
and the first characters of its \shaname{} digest, so a reader can check that the table and the file agree. No table or
figure shows an identifier, and no cell or plotted group describes fewer than $\proto{privacy.min_cell}$ customers.

\paragraph{Names.} Identifiers that contain digits, such as decision records, notebooks, rules and files, are set in
\ident{typewriter} type. Notebook sources are named by their folder and number, as in \ident{country\_template/13};
a notebook run once per market appears once. Markets are written ALL (the whole bank), AR (Argentina), CO (Colombia) and
MX (Mexico). The main copy is the dataset the pipelines serve; the backup is the second copy that arrived with it.

\paragraph{Notation.} Symbols are defined where they are first used and collected in \cref{app:notation}. Throughout,
$\E$, $\Prob$ and $\Var$ denote expectation, probability and variance, a hat marks an estimate, an asterisk a bootstrap
replicate, and $z_{1-\alpha/2}$ the standard normal quantile.

\paragraph{Rebuilding.} \Cref{app:software} gives the commands that rebuild the report and its numbers, the versions of
the software that computed them, the digests of the data files they were computed from, and the commits of the
protocol and of the build.
